\documentclass[10pt,a4paper]{amsart}
\usepackage{amsmath,amssymb,amsthm,graphicx}
% Retained minimal preamble for the Stacks Project chapter varieties.tex.
% Verbatim excerpt of the chapter preamble: Xy-pic setup (native preamble lines 14-18), the native theorem
% declarations with the mathematical macros (native preamble lines 154-217), and the project own no-op
% definitions of the reference/slogan/history display environments (native wrapper.tex lines 1-3) that the
% counted statement uses for its slogan. Copied mechanically; nothing added, removed or redefined.
% For commutative diagrams we use Xy-pic
\usepackage[all]{xy}

% We use 2cell for 2-commutative diagrams.
\xyoption{2cell}
% Theorem environments.
%
\theoremstyle{plain}
\newtheorem{theorem}[subsection]{Theorem}
\newtheorem{proposition}[subsection]{Proposition}
\newtheorem{lemma}[subsection]{Lemma}

\theoremstyle{definition}
\newtheorem{definition}[subsection]{Definition}
\newtheorem{example}[subsection]{Example}
\newtheorem{exercise}[subsection]{Exercise}
\newtheorem{situation}[subsection]{Situation}

\theoremstyle{remark}
\newtheorem{remark}[subsection]{Remark}
\newtheorem{remarks}[subsection]{Remarks}

\numberwithin{equation}{subsection}

% Macros
%
\def\lim{\mathop{\mathrm{lim}}\nolimits}
\def\colim{\mathop{\mathrm{colim}}\nolimits}
\def\Spec{\mathop{\mathrm{Spec}}}
\def\Hom{\mathop{\mathrm{Hom}}\nolimits}
\def\Ext{\mathop{\mathrm{Ext}}\nolimits}
\def\SheafHom{\mathop{\mathcal{H}\!\mathit{om}}\nolimits}
\def\SheafExt{\mathop{\mathcal{E}\!\mathit{xt}}\nolimits}
\def\Sch{\mathit{Sch}}
\def\Mor{\mathop{\mathrm{Mor}}\nolimits}
\def\Ob{\mathop{\mathrm{Ob}}\nolimits}
\def\Sh{\mathop{\mathit{Sh}}\nolimits}
\def\NL{\mathop{N\!L}\nolimits}
\def\CH{\mathop{\mathrm{CH}}\nolimits}
\def\proetale{{pro\text{-}\acute{e}tale}}
\def\etale{{\acute{e}tale}}
\def\QCoh{\mathit{QCoh}}
\def\Ker{\mathop{\mathrm{Ker}}}
\def\Im{\mathop{\mathrm{Im}}}
\def\Coker{\mathop{\mathrm{Coker}}}
\def\Coim{\mathop{\mathrm{Coim}}}

% Boxtimes
%
\DeclareMathSymbol{\boxtimes}{\mathbin}{AMSa}{"02}

%
% Macros for moduli stacks/spaces
%
\def\QCohstack{\mathcal{QC}\!\mathit{oh}}
\def\Cohstack{\mathcal{C}\!\mathit{oh}}
\def\Spacesstack{\mathcal{S}\!\mathit{paces}}
\def\Quotfunctor{\mathrm{Quot}}
\def\Hilbfunctor{\mathrm{Hilb}}
\def\Curvesstack{\mathcal{C}\!\mathit{urves}}
\def\Polarizedstack{\mathcal{P}\!\mathit{olarized}}
\def\Complexesstack{\mathcal{C}\!\mathit{omplexes}}
% \Pic is the operator that assigns to X its picard group, usage \Pic(X)
% \Picardstack_{X/B} denotes the Picard stack of X over B
% \Picardfunctor_{X/B} denotes the Picard functor of X over B
\def\Pic{\mathop{\mathrm{Pic}}\nolimits}
\def\Picardstack{\mathcal{P}\!\mathit{ic}}
\def\Picardfunctor{\mathrm{Pic}}
\def\Deformationcategory{\mathcal{D}\!\mathit{ef}}
\newenvironment{reference}{}{}
\newenvironment{slogan}{}{}
\newenvironment{history}{}{}
\usepackage{cleveref}
\crefname{originallemma}{original lemma}{original lemmas}
\Crefname{originallemma}{Original lemma}{Original lemmas}
\makeatletter
\newlabel{lemma-dimension-locally-algebraic}{{0A21}{source}{}{sourcepointer-1}{}}
\newlabel{lemma-dimension-locally-algebraic@cref}{{[originallemma][0][]0A21}{[1][0][]source}{}{}{}}
\makeatother
\crefname{originalreference}{original reference}{original references}
\Crefname{originalreference}{Original reference}{Original references}
\makeatletter
\newlabel{divisors-lemma-effective-Cartier-in-points}{{0AG8}{source}{}{sourcepointer-2}{}}
\newlabel{divisors-lemma-effective-Cartier-in-points@cref}{{[originalreference][0][]0AG8}{[1][0][]source}{}{}{}}
\makeatother
\crefname{originalreference}{original reference}{original references}
\Crefname{originalreference}{Original reference}{Original references}
\makeatletter
\newlabel{algebra-lemma-characterize-dvr}{{00PD}{source}{}{sourcepointer-3}{}}
\newlabel{algebra-lemma-characterize-dvr@cref}{{[originalreference][0][]00PD}{[1][0][]source}{}{}{}}
\makeatother
\crefname{originallemma}{original lemma}{original lemmas}
\Crefname{originallemma}{Original lemma}{Original lemmas}
\makeatletter
\newlabel{lemma-dense-smooth-open-variety-over-perfect-field}{{0B8X}{source}{}{sourcepointer-4}{}}
\newlabel{lemma-dense-smooth-open-variety-over-perfect-field@cref}{{[originallemma][0][]0B8X}{[1][0][]source}{}{}{}}
\makeatother
\crefname{originalreference}{original reference}{original references}
\Crefname{originalreference}{Original reference}{Original references}
\makeatletter
\newlabel{algebra-lemma-separable-smooth}{{00TV}{source}{}{sourcepointer-5}{}}
\newlabel{algebra-lemma-separable-smooth@cref}{{[originalreference][0][]00TV}{[1][0][]source}{}{}{}}
\makeatother
\title{Equivalence of regularity, normality, discrete valuation ring, invertible ideal sheaf and Cartier divisor conditions at a closed point of a curve}
\author{The Stacks Project Authors}
\date{stacks-project:varieties.tex:local-snapshot-20261011; GNU FDL 1.2 or later}
\begin{document}
\maketitle
\noindent\textit{Editorial note.} All delivered mathematical text is the author's own, copied verbatim from the pinned Stacks Project chapter file varieties.tex; no statement or proof step was rewritten, reordered or removed.\par\medskip
\noindent\textit{Editorial note.} The counted span cites five cross-chapter labels, each recorded in assembly.external\_references with its canonical tag as a navigation pointer.\par\medskip
\noindent\textit{Editorial note.} The proof cites no bibliography item inside the counted span.\par\medskip
\noindent\textit{Allowed standard prerequisites (editorial).}\begin{itemize}
\item Effective Cartier divisors in points and the characterization of discrete valuation rings. Stacks Project, Divisors, Lemma divisors-lemma-effective-Cartier-in-points (tag 0B3C); Algebra, Lemma algebra-lemma-characterize-dvr (tag 034Z). The first identifies the divisor conditions with principality of the maximal ideal and the second gives the equivalence with regularity and normality.
\item Density of the smooth locus and separability criteria for smoothness. Stacks Project, Varieties, Lemma lemma-dense-smooth-open-variety-over-perfect-field (tag 0B2J); Algebra, Lemma algebra-lemma-separable-smooth (tag 030V). Supply the smoothness equivalence under the perfect and separable hypotheses.
\end{itemize}\medskip
\ifdefined\hypertarget\hypertarget{sourcepointer-1}{}\fi
\noindent\textit{Source reference.} Original lemma 0A21: sources/stacks-chow/tags.txt, source lines 12689--12689. Cross-chapter standard prerequisite cited by the author inside the counted span; it is used as a pointer and its own proof is not part of this item.\par\medskip
\ifdefined\hypertarget\hypertarget{sourcepointer-2}{}\fi
\noindent\textit{Source reference.} Original reference 0AG8: sources/stacks-chow/tags.txt, source lines 13186--13186. Cross-chapter standard prerequisite cited by the author inside the counted span; it is used as a pointer and its own proof is not part of this item.\par\medskip
\ifdefined\hypertarget\hypertarget{sourcepointer-3}{}\fi
\noindent\textit{Source reference.} Original reference 00PD: sources/stacks-chow/tags.txt, source lines 927--927. Cross-chapter standard prerequisite cited by the author inside the counted span; it is used as a pointer and its own proof is not part of this item.\par\medskip
\ifdefined\hypertarget\hypertarget{sourcepointer-4}{}\fi
\noindent\textit{Source reference.} Original lemma 0B8X: sources/stacks-chow/tags.txt, source lines 14155--14155. Cross-chapter standard prerequisite cited by the author inside the counted span; it is used as a pointer and its own proof is not part of this item.\par\medskip
\ifdefined\hypertarget\hypertarget{sourcepointer-5}{}\fi
\noindent\textit{Source reference.} Original reference 00TV: sources/stacks-chow/tags.txt, source lines 1089--1089. Cross-chapter standard prerequisite cited by the author inside the counted span; it is used as a pointer and its own proof is not part of this item.\par\medskip
\begin{lemma}
\label{lemma-regular-point-on-curve}
Let $k$ be a field. Let $X$ be a curve over $k$. Let $x \in X$ be a closed
point. We think of $x$ as a (reduced) closed subscheme of $X$ with sheaf
of ideals $\mathcal{I}$. The following are equivalent
\begin{enumerate}
\item $\mathcal{O}_{X, x}$ is regular,
\item $\mathcal{O}_{X, x}$ is normal,
\item $\mathcal{O}_{X, x}$ is a discrete valuation ring,
\item $\mathcal{I}$ is an invertible $\mathcal{O}_X$-module,
\item $x$ is an effective Cartier divisor on $X$.
\end{enumerate}
If $k$ is perfect or if $\kappa(x)$ is separable over $k$,
these are also equivalent to
\begin{enumerate}
\item[(6)] $X \to \Spec(k)$ is smooth at $x$.
\end{enumerate}
\end{lemma}

\begin{proof}
Since $X$ is a curve, the local ring $\mathcal{O}_{X, x}$ is a Noetherian
local domain of dimension $1$ (Lemma \ref{lemma-dimension-locally-algebraic}).
Parts (4) and (5) are equivalent by definition and are equivalent to
$\mathcal{I}_x = \mathfrak m_x \subset \mathcal{O}_{X, x}$ having one generator
(Divisors, Lemma \ref{divisors-lemma-effective-Cartier-in-points}).
The equivalence of (1), (2), (3), (4), and (5) therefore follows from
Algebra, Lemma \ref{algebra-lemma-characterize-dvr}. The final statement
follows from Lemma \ref{lemma-dense-smooth-open-variety-over-perfect-field}
in case $k$ is perfect. If $\kappa(x)/k$ is separable, then the equivalence
follows from Algebra, Lemma \ref{algebra-lemma-separable-smooth}.
\end{proof}

\end{document}
